% GENERATED FILE. Edit canonical lessons, then run npm run generate:books.
% canonical-source-hash: fnv1a64:5bbe97192e85018b
% canonical-lessons: HI-C441-savera, HI-C441-suryoday, HI-C441-suryast, HI-C441-ksitij, HI-C441-godhuli

\chapter{Dawn and Dusk}
\label{ch:hi-dawn-and-dusk}
\hlchaptermodality{441}

\begin{chapteropening}
\textbf{By the end of this chapter:} \emph{I can talk about early morning, sunrise and sunset, the horizon and dusk.}

Complete the last lesson of chapter 441: I can talk about early morning, sunrise and sunset, the horizon and dusk.
\end{chapteropening}

\section[\texorpdfstring{\hi{सवेरा} (saverā)}{सवेरा (saverā)}]{\hi{सवेरा} (saverā) — early morning, daybreak}

\label{lesson:HI-C441-savera}



\begin{quote}
Before the new one: say the Hindi for to clear (of clouds), then the Hindi for to blow away.
\end{quote}

\subsection*{You'll want to know: \hi{सवेरा}}
\textbf{\hi{सवेरा}} — \emph{saverā} — \textquotedblleft{}early morning, daybreak\textquotedblright{}.

\begin{quote}
\hi{हम} \hi{सवेरे} \hi{पाँच} \hi{बजे} \hi{निकले।}
\end{quote}

\begin{quote}
\emph{ham savere pā̃c baje nikle.}
\end{quote}

\begin{quote}
We set off at five in the early morning.
\end{quote}

Masculine.

\begin{cousinweb}
\textbf{\hi{सवेरा}} (\emph{saverā}) goes back to \emph{*savēla}, \textquotedblleft{}in time, early\textquotedblright{}, a form CDIAL reconstructs (CDIAL 13291), which lists Hindi \emph{saverā}, \textquotedblleft{}early morning\textquotedblright{}.
\end{cousinweb}

\begin{grammarlens}[title={What you've built}]
One more word for getting around.
\end{grammarlens}

\subsection*{Your turn}
\begin{itemize}
\raggedright
  \item \emph{Say it:} \emph{saverā}
  \item \emph{Say it:} \emph{saverā}, once more
  \item \emph{Say it:} \emph{ham savere pā̃c baje nikle.} — We set off at five in the early morning.
  \item \emph{Recall:} say the Hindi for to clear (of clouds), then the Hindi for to blow away, then say \emph{saverā} again
\end{itemize}

\begin{tcolorbox}[breakable,colback=teal!4,colframe=teal!35!black,title={Before you move on}]
What does \hi{सवेरा} mean? (\textquotedblleft{}Early morning, daybreak\textquotedblright{}.) Say it once more.
\end{tcolorbox}
\section[\texorpdfstring{\hi{सूर्योदय} (sūryoday)}{सूर्योदय (sūryoday)}]{\hi{सूर्योदय} (sūryoday) — sunrise}

\label{lesson:HI-C441-suryoday}



\begin{quote}
Before the new one: say the Hindi for to blow away, then the Hindi for daybreak.
\end{quote}

\subsection*{You'll want to know: \hi{सूर्योदय}}
\textbf{\hi{सूर्योदय}} — \emph{sūryoday} — \textquotedblleft{}sunrise\textquotedblright{}.

\begin{quote}
\hi{पहाड़} \hi{से} \hi{सूर्योदय} \hi{बहुत} \hi{सुंदर} \hi{दिखता} \hi{है।}
\end{quote}

\begin{quote}
\emph{pahāṛ se sūryoday bahut sundar dikhtā hai.}
\end{quote}

\begin{quote}
The sunrise looks beautiful from the mountain.
\end{quote}

Masculine.

\begin{cousinweb}
\textbf{\hi{सूर्योदय}} (\emph{sūryoday}) is a tatsama, Sanskrit \emph{sūryodaya} (Monier-Williams), from \emph{sūrya}, \textquotedblleft{}the sun\textquotedblright{}, and \emph{udaya}, \textquotedblleft{}going up\textquotedblright{}; the two vowels a + u have merged into o.
\end{cousinweb}

\begin{grammarlens}[title={What you've built}]
One more word for getting around.
\end{grammarlens}

\subsection*{Your turn}
\begin{itemize}
\raggedright
  \item \emph{Say it:} \emph{sūryoday}
  \item \emph{Say it:} \emph{sūryoday}, once more
  \item \emph{Say it:} \emph{pahāṛ se sūryoday bahut sundar dikhtā hai.} — The sunrise looks beautiful from the mountain.
  \item \emph{Recall:} say the Hindi for to blow away, then the Hindi for daybreak, then say \emph{sūryoday} again
\end{itemize}

\begin{tcolorbox}[breakable,colback=teal!4,colframe=teal!35!black,title={Before you move on}]
What does \hi{सूर्योदय} mean? (\textquotedblleft{}Sunrise\textquotedblright{}.) Say it once more.
\end{tcolorbox}
\section[\texorpdfstring{\hi{सूर्यास्त} (sūryāst)}{सूर्यास्त (sūryāst)}]{\hi{सूर्यास्त} (sūryāst) — sunset}

\label{lesson:HI-C441-suryast}



\begin{quote}
Before the new one: say the Hindi for daybreak, then the Hindi for sunrise.
\end{quote}

\subsection*{You'll want to know: \hi{सूर्यास्त}}
\textbf{\hi{सूर्यास्त}} — \emph{sūryāst} — \textquotedblleft{}sunset\textquotedblright{}.

\begin{quote}
\hi{हम} \hi{समुद्र} \hi{के} \hi{किनारे} \hi{सूर्यास्त} \hi{देख} \hi{रहे} \hi{थे।}
\end{quote}

\begin{quote}
\emph{ham samudra ke kināre sūryāst dekh rahe the.}
\end{quote}

\begin{quote}
We were watching the sunset by the sea.
\end{quote}

Masculine.

\begin{cousinweb}
\textbf{\hi{सूर्यास्त}} (\emph{sūryāst}) is a tatsama, Sanskrit \emph{sūryāsta}, \textquotedblleft{}sunset\textquotedblright{} (Monier-Williams), from \emph{sūrya}, \textquotedblleft{}the sun\textquotedblright{}, and \emph{asta}, \textquotedblleft{}setting\textquotedblright{}.
\end{cousinweb}

\begin{grammarlens}[title={What you've built}]
One more word for getting around.
\end{grammarlens}

\subsection*{Your turn}
\begin{itemize}
\raggedright
  \item \emph{Say it:} \emph{sūryāst}
  \item \emph{Say it:} \emph{sūryāst}, once more
  \item \emph{Say it:} \emph{ham samudra ke kināre sūryāst dekh rahe the.} — We were watching the sunset by the sea.
  \item \emph{Recall:} say the Hindi for daybreak, then the Hindi for sunrise, then say \emph{sūryāst} again
\end{itemize}

\begin{tcolorbox}[breakable,colback=teal!4,colframe=teal!35!black,title={Before you move on}]
What does \hi{सूर्यास्त} mean? (\textquotedblleft{}Sunset\textquotedblright{}.) Say it once more.
\end{tcolorbox}
\section[\texorpdfstring{\hi{क्षितिज} (kṣitij)}{क्षितिज (kṣitij)}]{\hi{क्षितिज} (kṣitij) — the horizon}

\label{lesson:HI-C441-ksitij}



\begin{quote}
Before the new one: say the Hindi for sunrise, then the Hindi for sunset.
\end{quote}

\subsection*{You'll want to know: \hi{क्षितिज}}
\textbf{\hi{क्षितिज}} — \emph{kṣitij} — \textquotedblleft{}the horizon\textquotedblright{}.

\begin{quote}
\hi{सूरज} \hi{क्षितिज} \hi{के} \hi{पीछे} \hi{छिप} \hi{गया।}
\end{quote}

\begin{quote}
\emph{sūraj kṣitij ke pīche chip gayā.}
\end{quote}

\begin{quote}
The sun went down behind the horizon.
\end{quote}

Masculine.

\begin{cousinweb}
\textbf{\hi{क्षितिज}} (\emph{kṣitij}) is a tatsama, Sanskrit \emph{kṣitija}, \textquotedblleft{}born of the earth\textquotedblright{}, which the astronomers used for the horizon (Monier-Williams): where the sky seems to rise out of the earth.
\end{cousinweb}

\begin{grammarlens}[title={What you've built}]
One more word for getting around.
\end{grammarlens}

\subsection*{Your turn}
\begin{itemize}
\raggedright
  \item \emph{Say it:} \emph{kṣitij}
  \item \emph{Say it:} \emph{kṣitij}, once more
  \item \emph{Say it:} \emph{sūraj kṣitij ke pīche chip gayā.} — The sun went down behind the horizon.
  \item \emph{Recall:} say the Hindi for sunrise, then the Hindi for sunset, then say \emph{kṣitij} again
\end{itemize}

\begin{tcolorbox}[breakable,colback=teal!4,colframe=teal!35!black,title={Before you move on}]
What does \hi{क्षितिज} mean? (\textquotedblleft{}The horizon\textquotedblright{}.) Say it once more.
\end{tcolorbox}
\section[\texorpdfstring{\hi{गोधूलि} (godhūli)}{गोधूलि (godhūli)}]{\hi{गोधूलि} (godhūli) — dusk (when the cattle come home)}

\label{lesson:HI-C441-godhuli}



\begin{quote}
Before the new one: say the Hindi for sunset, then the Hindi for the horizon.
\end{quote}

\subsection*{You'll want to know: \hi{गोधूलि}}
\textbf{\hi{गोधूलि}} — \emph{godhūli} — \textquotedblleft{}dusk (when the cattle come home)\textquotedblright{}.

\begin{quote}
\hi{गोधूलि} \hi{में} \hi{गायें} \hi{गाँव} \hi{लौट} \hi{रही} \hi{थीं।}
\end{quote}

\begin{quote}
\emph{godhūli meṁ gāyẽ gā̃v lauṭ rahī thīṁ.}
\end{quote}

\begin{quote}
At dusk the cows were coming back to the village.
\end{quote}

Feminine.

\begin{cousinweb}
\textbf{\hi{गोधूलि}} (\emph{godhūli}) is a tatsama, Sanskrit \emph{godhūli} (Monier-Williams), literally the dust of \emph{go}, a word for both the cow and the earth. Monier-Williams reads it as earth-dust, the haze that seems to rise at that hour; the usual telling is the dust raised by the cattle coming home. It is still an auspicious time for weddings.
\end{cousinweb}

\begin{grammarlens}[title={What you've built}]
One more word for getting around.
\end{grammarlens}

\subsection*{Your turn}
\begin{itemize}
\raggedright
  \item \emph{Say it:} \emph{godhūli}
  \item \emph{Say it:} \emph{godhūli}, once more
  \item \emph{Say it:} \emph{godhūli meṁ gāyẽ gā̃v lauṭ rahī thīṁ.} — At dusk the cows were coming back to the village.
  \item \emph{Recall:} say the Hindi for sunset, then the Hindi for the horizon, then say \emph{godhūli} again
\end{itemize}

\begin{tcolorbox}[breakable,colback=teal!4,colframe=teal!35!black,title={Before you move on}]
What does \hi{गोधूलि} mean? (\textquotedblleft{}Dusk (when the cattle come home)\textquotedblright{}.) Say it once more.
\end{tcolorbox}
